\begin{claim}{C.63}{Proved}{A controlled moving-base lift on principal orbits with fixed frontier}
Let $f_t$ be the original based ambient-null plane loop in
the centralizer of an arbitrary operator group $H$, and let
$G$ be the group generated by its values. Let $W$ be an
$H$-invariant union of open $G$-orbits. For every $O\subset W$
suppose that the full local centralizer
\[
 C_O=C_{\operatorname{Homeo}(O)}(G|_O)
\]
has dimension one, with its actual compact-open topology.
Suppose also that every point of $\partial W$ is fixed by $G$.
Then the actual global loop
\[
 \widehat f_t(x)=
 \begin{cases}f_t(x),&x\in W,\\x,&x\notin W\end{cases}
                                                               \tag{C.63.1}
\]
has a based null-homotopy in the $H$-centralizer, fixing the
complement of $W$ throughout. Both this homotopy and its
inverse are jointly continuous across every accumulation
of open orbits.

On each $O$, the induced contraction of the actual principal
quotient $O/C_O^0$ is the median contraction of its original
based real target lift. The original source cover, every
deck label, and the exact $H_O$ restriction kernel are
retained. No nontrivial global quotient to a nonopen orbit
is required, so the assertion includes the empty global
target family. In particular it contracts the original
loop when $\R^2\setminus W$ is fixed by $G$.
\end{claim}

\begin{proof}
\emph{The actual principal coordinates and source labels.}
C.30 and C.33 make $C_O$ a free proper Lie transformation
group. Put $A=C_O^0$. Their actual principal bundle and
section give
\[
 \Theta:A\times B\xrightarrow{\ \cong\ }O,
 \qquad\Theta(a,b)=a\,s(b),\qquad B=O/A,
                                                               \tag{C.63.2}
\]
where each of $A,B$ is a line or a circle. The pair of
circles is excluded because $O$ is noncompact.
Let $\eta_A:\R\to A$ and $\eta_B:\R\to B$ be the identity
on a line factor and the specified quotient by $\Z$ on
a circle factor. The map
\[
               p_*(u,v)=\Theta(\eta_A(u),\eta_B(v))
                                                               \tag{C.63.3}
\]
is a universal covering. With basepoints fixed, its unique
covering isomorphism with the original $p_O:X_O\to O$ is
a homeomorphism $\chi_O:X_O\to\R^2$ satisfying
$p_O=p_*\chi_O$. The inverse is the based lift in the
opposite direction; uniqueness of based lifts proves
both inverse identities. Hence
\[
 F_O=\operatorname{pr}_2\chi_O,
 \qquad
 \chi_O\delta\chi_O^{-1}(u,v)
       =(u+m_A(\delta),v+m_B(\delta))
                                                               \tag{C.63.4}
\]
for each original deck map $\delta$. A line factor has
zero label, and a circle factor has the original integer
label under this based covering comparison. At most one
factor is a circle.

Every element of $C_O$ normalizes its identity component
$A$ and preserves its orbit partition. C.33 proves that
$C_O/A$ acts freely and properly on $B$, so no induced
base map is decreasing: a decreasing line or circle
homeomorphism has a fixed point. Retain the complete
affine fiber changes constructed in C.33. On the original
universal cover, the full group of lifts of $C_O$ is
generated by all fiber translations and, when its real
base image is nontrivial, one map of the form
\[
 T_c(u,v)=(u+c,v),\qquad
 q(u,v)=(\lambda u+\kappa,v+d),
 \quad d>0,\quad\lambda\ne0.                            \tag{C.63.5}
\]
This also includes the source deck group. For a circle
base its lift labels supply the base period; for a
circle fiber its integer deck maps are among the
$T_c$. The constant $\kappa$ retains the choice of
component representative and lift. C.33's fiber gauge
removes the nonconstant translation term; a different
original lift only changes the constant term. The
negative values of $\lambda$ remain allowed. If the
base image is trivial, only $T_c$ occurs.

The based source lift of the original loop commutes
with every original lift of an operator in $C_O$.
Indeed its commutator with such a lift is deck-valued
and is initially the identity. Commutation with all
$T_c$ gives the exact expression
\[
 \widetilde f_t^O(u,v)=(u+\phi_t^O(v),\beta_t^O(v)),
                                                               \tag{C.63.6}
\]
where $\beta_t^O$ is increasing. Here and below this
formula is written through $\chi_O$. By C.32 the
ambient-null loop restricts to an ambient-null loop
on $O$, so its based source lift closes. Thus
\[
 \phi_0^O=\phi_1^O=0,
       \qquad\beta_0^O=\beta_1^O=\id.                 \tag{C.63.7}
\]
When (C.63.5) has the additional generator, its exact
commutation equations are
\[
 \phi_t^O(v+d)=\lambda\phi_t^O(v),\qquad
 \beta_t^O(v+d)=\beta_t^O(v)+d.                         \tag{C.63.8}
\]
The constant $\kappa$ cancels from this comparison,
rather than being discarded from the operator group.

The base maps $\beta_t^O$ have no common fixed point.
A common fixed point would project to a point of $B$
fixed by every original generator, contradicting the
transitivity of $G$ on $O$ and hence on $B$.
For any $a\in\R$ and any $\delta>0$, the compact
increment family
\[
 \{\beta_s^O(\beta_t^O)^{-1}:|s-t|\leq\delta\}
                                                               \tag{C.63.9}
\]
sends $a$ to points on both sides of $a$. To check
this for arbitrarily short increments, the path
$t\mapsto(\beta_t^O)^{-1}(a)$ is nonconstant. It
therefore has unequal values at arbitrarily close
times $s,t$; otherwise it is locally constant and
then constant on the connected parameter interval.
The corresponding increment moves $a$, and its
inverse moves $a$ to the opposite side. The
parameter set $|s-t|\leq\delta$ is convex, contains
the diagonal, and gives a path-connected increment
family. Its evaluation image consequently contains
a neighborhood of $a$. This uses inverse evaluation
at the specified point, not a displacement at an
unrelated moving point.

\emph{One compact family of global maps and original lifts.}
There are at most countably many open orbits, since
each contains a member of a fixed countable basis
disjoint from all the others. Consider the product
of the global plane homeomorphism group and the
homeomorphism groups of all the selected original
source covers. In this product take the closed
subgroup $\mathcal E$ of compatible tuples
\[
             (k,(\widetilde k^O)_{O\subset W}),
       \qquad p_O\widetilde k^O=k p_O,                 \tag{C.63.10}
\]
which commute globally with $H$ and in each cover
with every retained lift of $C_O$.
All these conditions are closed pointwise
equalities, including the inverse equalities.
Compatibility implies $k(O)=O$, using the
surjectivity of $p_O$ and the inverse tuple.
Alternatively each original word preserves the
closed complement of every $O$; limits of both
maps and inverses preserve that complement too.
Thus no limiting homeomorphism sends an open
orbit into its frontier.

Each factor is Polish in its compact-open
homeomorphism topology, with inverses controlled
on compact sets. The countable product and its
closed subgroup $\mathcal E$ are Polish. Choose
a complete compatible metric on $\mathcal E$
controlling these maps and inverses. For example,
after choosing a complete compatible metric,
add its value on the inverse pair; the resulting
metric is still complete and compatible.
The original path
\[
             L_t=(f_t,(\widetilde f_t^O)_O)
                                                               \tag{C.63.11}
\]
is a continuous based loop in this exact group.

We construct positive numbers $\delta_n$ so that
all the products
\[
 L_{s_n}L_{t_n}^{-1}\cdots
 L_{s_1}L_{t_1}^{-1}L_r,
       \qquad |s_j-t_j|\leq\delta_j,\quad r\in[0,1],
                                                               \tag{C.63.12}
\]
converge uniformly over their entire parameter
product as $n\to\infty$. Suppose the first $n-1$
choices have been made; their product image is
compact. Continuity of left multiplication,
uniform on this compact image, gives an identity
neighborhood whose multiplication changes every
such product by less than $2^{-n}$ in the chosen
metric. Uniform continuity of
$(s,t)\mapsto L_sL_t^{-1}$ near its compact
diagonal gives $\delta_n$ placing all the next
increments in that neighborhood. The successive
product maps are therefore uniformly Cauchy.
Completeness gives their uniform limit in
$\mathcal E$, including every global inverse
and every original lifted inverse.

Put
\[
 T=[0,1]\times\prod_{n\geq1}
       \{(s_n,t_n)\in[0,1]^2:|s_n-t_n|\leq\delta_n\}.
                                                               \tag{C.63.13}
\]
This is compact and path connected; coordinatewise
straight segments give paths with their exact
labels. The uniform limit just constructed is a
continuous map $T\to\mathcal E$. Denote its
labelled image by $\mathcal K$. Let $\mathcal K_n$
be the family of products through the first $n$
increments, with $\mathcal K_0=\{L_r\}$.
Setting all later increments to the identity
shows $\mathcal K_n\subset\mathcal K$ and
$\mathcal K_n\subset\mathcal K_{n+1}$.
Every global value of $\mathcal K$ commutes with
$H$ and fixes every point fixed by $G$.
The source labels in $\mathcal K$ are part of
the construction even when distinct parameters
have the same global homeomorphism.

The C.62 triangle works for any path in such a
labelled family. Its proof uses the relative
maps $h_r=k_rk_a^{-1}$, so it does not require
the selected path itself to start at the identity.
Connect any two chosen labels by a path in $T$.
Their equal-real-level fiber interval is contained
in the corresponding image of
\[
             \mathcal K\mathcal K^{-1}\mathcal K.
                                                               \tag{C.63.14}
\]
The projected family in (C.63.14) is compact in
the global plane homeomorphism group. The same
argument for the inverse labelled family uses
$\mathcal K^{-1}\mathcal K\mathcal K^{-1}$.
No continuous choice of connecting paths or
endpoint labels is needed for these uniform
containments.

\emph{Enlarging the allowed fiber intervals.}
Fix one source orbit and suppress $O$ from the
notation. Let $B_{\rho,t}$ be C.27's median of
the real path $\beta_t$, using prefix length
$\rho$. Thus $B_{0,t}=\id$, $B_{1,t}=\beta_t$,
and $B_{\rho,0}=B_{\rho,1}=\id$.
For every $v$, its value lies between $v$ and
$\beta_t(v)$, since every running minimum and
maximum contains the initial value $v$ between
them. Hence some original prefix time between
$0$ and $t$ attains $B_{\rho,t}(v)$.

Each labelled source map $k\in\mathcal K$ has
the form
\[
                    k(u,v)=(u+c_k(v),b_k(v)),
                                                               \tag{C.63.15}
\]
with its jointly continuous real functions in
the label and $v$. For $p=(\rho,t,v)$ define
the nonempty compact interval
\[
 A_n(p)=\operatorname{conv}
       \{c_k(v):k\in\mathcal K_n,
                          \ b_k(v)=B_{\rho,t}(v)\}.
                                                               \tag{C.63.16}
\]
Nonemptiness follows from the preceding prefix
observation, and compactness follows from the
closed equality in the compact labelled
parameter set. The identical definition with
$\mathcal K$ gives a compact interval
$A_\infty(p)$ containing every $A_n(p)$.

Here is the precise lower-limit property used
for selection:
\[
             A_n(p)\subseteq
                    \liminf_{p'\to p}A_{n+1}(p').       \tag{C.63.17}
\]
First take an endpoint value $c_k(v)$ before
convexification, with its fixed label
$k\in\mathcal K_n$ and $b_k(v)=a=B_{\rho,t}(v)$.
At nearby parameters $p'$, both $b_k(v')$ and
$B_{\rho',t'}(v')$ tend to $a$. By (C.63.9),
increments with arbitrarily small time
differences move $a$ to both sides of itself.
Choose two such increments within the
$(n+1)$st allowed family, as close to the
identity as desired. Their strict inequalities
continue to bracket the desired target at
all sufficiently nearby $p'$. A path in that
small increment family and the intermediate
value theorem therefore give an increment
$e_{p'}$ with
\[
          b_{e_{p'}k}(v')=B_{\rho',t'}(v').
                                                               \tag{C.63.18}
\]
As the chosen increment neighborhood decreases
to the identity, its effect on the fiber
coordinate also tends uniformly to zero near
the specified source point. Thus the value
$c_k(v)$ is approximated by allowed values
of $\mathcal K_{n+1}$ at every nearby parameter.
For a point between two allowed endpoint
values, perform this construction for both
endpoints and take the same convex combination.
This proves (C.63.17). It does not assert
lower semicontinuity of $A_n$ itself.

The operator equations are retained by doing
the selection in their actual line bundle.
If (C.63.5) has a base generator, all members
of $\mathcal K$ satisfy
\[
 c_k(v+d)=\lambda c_k(v),\qquad
 b_k(v+d)=b_k(v)+d.                                    \tag{C.63.19}
\]
The median has the same base equation.
Thus the intervals (C.63.16) define convex
fiber sets in the real line bundle
\[
 (v,z)\sim(v+d,\lambda z)
       \quad\hbox{over }\R/d\Z,                       \tag{C.63.20}
\]
pulled back over the parameter square.
Its continuous fiber norm is
$\|[v,z]\|=|\lambda|^{-v/d}|z|$.
For trivial base image use the ordinary real
line bundle over $[0,1]^2\times\R$.
Both bases are metrizable and paracompact.
Negative $\lambda$ gives the corresponding
nonorientable line bundle and preserves its
convex fiber sets just as in (C.63.19).

\emph{The relative continuous-selection construction.}
On the closed union $D$ of the four boundary
faces of the parameter square prescribe
\[
 c(0,t,v)=0,\quad c(1,t,v)=\phi_t(v),\quad
 c(\rho,0,v)=c(\rho,1,v)=0.                            \tag{C.63.21}
\]
This is a continuous section on $D$, including
the corners by (C.63.7), and belongs to every
allowed interval there. Replace $A_n$ by this
prescribed singleton on $D$ and leave it
unchanged elsewhere; denote the new fiber
sets by $F_n$. Property (C.63.17) survives.
At a point of $D$ the required inclusion for
approaching points off $D$ is the already
proved inclusion for its prescribed value,
and for approaching points in $D$ it is
continuity of (C.63.21). Away from $D$ a
neighborhood avoids $D$.

We give the elementary approximation proof
needed for these interlaced intervals. It
also applies in the line bundle (C.63.20).
Choose positive $\epsilon_n=2^{-n}$.
At each parameter point choose one value
of $F_0$ and extend it as a constant local
section in a bundle chart. By the lower-limit
inclusion, a sufficiently small neighborhood
makes this local section have distance less
than $\epsilon_1$ from $F_1$.
A locally finite partition of unity subordinate
to these neighborhoods combines the sections
to a continuous section $a_1$ with that same
distance bound. Indeed, at a given point the
corresponding nearby elements of $F_1$ have
their convex combination in $F_1$, so the
norm of the error is bounded by the same
weighted error bound.

Suppose $a_n$ has distance less than
$\epsilon_n$ from $F_n$. At each parameter
point choose a value in $F_n$ that close to
$a_n$, and extend it in a local bundle chart.
Shrink its neighborhood so that the extended
section remains within $2\epsilon_n$ of
$a_n$ and within $\epsilon_{n+1}$ of
$F_{n+1}$; the latter requirement is exactly
the lower-limit inclusion. Another locally
finite partition of unity gives $a_{n+1}$
with
\[
 \|a_{n+1}-a_n\|<2\epsilon_n,
       \qquad \operatorname{dist}(a_{n+1},F_{n+1})
                                      <\epsilon_{n+1}.
                                                               \tag{C.63.22}
\]
The summable first bounds give a limit section
$c$, uniformly in the fiber norm and hence
locally uniformly in every bundle chart.
It is continuous. At each fixed parameter,
every $F_n$ lies in the closed compact interval
$A_\infty$, and the second bound tends to zero;
therefore $c\in A_\infty$ at that parameter.
On $D$ the sets $F_n$ are the same prescribed
singleton, so $c$ equals (C.63.21) exactly.
This proves the required relative selection
without assuming that the original prefix
hulls themselves admit a continuous selection.

\emph{The intrinsic homeomorphisms and all operator equations.}
In the original real source chart set
\[
 K_{\rho,t}^O(u,v)
       =(u+c(\rho,t,v),B_{\rho,t}(v)).                 \tag{C.63.23}
\]
This is a jointly continuous family of
homeomorphisms, with inverse
\[
 (u,w)\longmapsto
 \bigl(u-c(\rho,t,B_{\rho,t}^{-1}(w)),
                         B_{\rho,t}^{-1}(w)\bigr).
                                                               \tag{C.63.24}
\]
The inverse family is jointly continuous
because the base median has a jointly
continuous inverse. The selected section
of (C.63.20) satisfies the first equation
of (C.63.19), and the base median satisfies
the second. Equations (C.63.5) therefore
commute exactly with (C.63.23); all fiber
translations commute as well. Consequently
every original source deck map commutes,
and (C.63.23) descends through $p_O$.
It commutes with the entire $C_O$ downstairs.
Equations (C.63.21) give the based
null-homotopy endpoint conditions.

Choose this construction on one representative
of each $H$-orbit of open components in $W$
and transport it by conjugation. The exact
restriction $H_O\to C_O$ has the kernel
specified in C.30.4. Commutation with its
image makes transport independent of the
chosen representative operator; its kernel
acts identically on $O$. Conjugation carries
$C_O^0$ to $C_{hO}^0$ and hence transports
the actual principal fibers. Original based
source lifts intertwine under a lift of $h$
by uniqueness at time zero. Their finite
products and the labelled limits above
intertwine too. Thus the original covering
labels are retained under this transport.

\emph{The compact-family bounds and fixed-frontier pasting.}
At every source point, membership
$c(\rho,t,v)\in A_\infty(\rho,t,v)$ places
the value (C.63.23) in the real fiber interval
between two labelled $\mathcal K$-values at
that same target level. The endpoints exist
because the defining level set in the compact
label space is compact. By (C.63.14), after
projection the value belongs to
\[
                   \mathcal K\mathcal K^{-1}\mathcal K(x).
                                                               \tag{C.63.25}
\]
For the inverse, put $y=(K_{\rho,t}^O)^{-1}x$.
The point $x$ is between two labelled values
$k_1y,k_2y$ in one real fiber. Since both
fiber maps are increasing, $y$ lies between
$k_1^{-1}x$ and $k_2^{-1}x$ in their common
real input fiber. The inverse version of
(C.63.14) therefore gives
\[
        (K_{\rho,t}^O)^{-1}(x)
             \in\mathcal K^{-1}\mathcal K\mathcal K^{-1}(x).
                                                               \tag{C.63.26}
\]
These are inclusions in one compact family
of actual global maps, independent of the
source orbit, the selected section, and all
homotopy parameters. They survive $H$
transport because every global member of
$\mathcal K$ commutes with $H$.

Define the global $K_{\rho,t}$ by these
transported maps on $W$ and the identity
outside $W$. It is a bijection: each open
orbit is preserved and the complement is
fixed. Let $x_n\to z\in\partial W$ and
allow the source orbit and both parameters
to vary arbitrarily. Every global map in
either compact family of (C.63.25)--(C.63.26)
fixes $z$, because $z$ is $G$-fixed.
Joint continuity of the compact global
evaluation family makes all its values at
$x_n$ tend uniformly to $z$.
The two inclusions thus prove continuity
of both the pasted map and its inverse at
every such limit. Interior continuity was
already proved, and points of the exterior
are fixed. Hence the global family consists
of homeomorphisms and is jointly continuous
with its inverse. It commutes with $H$ by
the transported construction and the
$H$-invariance of $W$. Its four endpoint
faces are precisely those required in
(C.63.1).

This establishes an actual moving-base lift
and global contraction at fixed frontiers,
including arbitrarily many open components
and their accumulation. At a nonfixed
frontier orbit, a compact global family need
not have one evaluation value; the argument
does not select the prescribed contraction
there. Open components with zero-dimensional
full local centralizer are also outside the
stated principal-fiber assertion. These
remaining cases of C.12 are retained.
\end{proof}
